% answers_short.tex —— 2012 年·福建卷·文科 · 简短答案速查
% 由 tools/make_short_answers.py 自 answers.tex 的 \gksolhead/\gkanswer 提取，
% 并与 2012/2012_数学真题_福建文_答案要点.md 逐题比对后生成，勿手改（改源 MD 后重跑）。

\begin{gknotice}
\textbf{使用说明}：本节只给最终答案便于快速核对，完整推导见解析版。答案为编写者独立推导，\textbf{非官方标准答案}。
\end{gknotice}

\gkpart{选择题}
\begin{center}\begin{tabular}{ccccccccccccc}
\toprule
题号 & 1 & 2 & 3 & 4 & 5 & 6 & 7 & 8 & 9 & 10 & 11 & 12 \\
\midrule
答案 & A & D & D & D & C & A & B & C & B & B & A & C \\
\bottomrule
\end{tabular}\end{center}

\gkpart{填空题}
\begin{itemize}[leftmargin=2.6em]
  \item[\textbf{13.}] $\sqrt2$。由正弦定理 $AC=\dfrac{BC\sin45^\circ}{\sin60^\circ}=\dfrac{\sqrt3\cdot\frac{\sqrt2}2}{\frac{\sqrt3}2}=\sqrt2$。
  \item[\textbf{14.}] $12$。女运动员 $98-56=42$ 人，按比例应抽 $28\times\dfrac{42}{98}=12$ 人。
  \item[\textbf{15.}] $(0,8)$。二次项系数为正，恒成立当且仅当判别式 $a^2-8a<0$，解得 $0<a<8$。
  \item[\textbf{16.}] $16$。按费用从小到大取边，依次取 $FG=1$、$AE=2$、$GD=2$、$AF=3$、$GC=3$、$BC=5$，六条边连通七个城市且不成圈，总费用 $16$。
\end{itemize}

\gkpart{解答题}
\begin{itemize}[leftmargin=2.6em]
  \item[\textbf{17.}] (1) $a_n=n$，$b_n=2^{n-1}$。由 $S_{10}=5(2+9d)=55$ 得 $d=1$；由 $b_1q^3=8$ 得 $q=2$；(2) $\dfrac29$。基本事件共 $3\times3=9$ 个，两项相等的情形为 $(1,1)$、$(2,2)$，共 $2$ 个
  \item[\textbf{18.}] (1) $\hat y=-20x+250$。$\bar x=8.5$，$\bar y=80$，故 $a=\bar y-b\bar x=80+170=250$；(2) $8.25$ 元。利润 $L=(x-4)(-20x+250)=-20x^2+330x-1000$，在 $x=\dfrac{330}{40}=8.25$ 处取最大值
  \item[\textbf{19.}] (1) $\dfrac13$。$\triangle MCC_1$ 的面积为 $\dfrac12\times2\times1=1$，$A$ 到平面 $MCC_1$ 的距离为 $1$，故 $V=\dfrac13\times1\times1=\dfrac13$；(2) 见解析版
  \item[\textbf{20.}] (1) $\dfrac34$。取②：$\sin^215^\circ+\cos^215^\circ-\sin15^\circ\cos15^\circ=1-\dfrac12\sin30^\circ=\dfrac34$；(2) 见解析版
  \item[\textbf{21.}] (1) $x^2=4y$。由对称性设 $A(-4\sqrt3,h)$，$B(4\sqrt3,h)$，由 $|OA|=8\sqrt3$ 得 $h=12$，代入 $48=2p\cdot12$ 得 $p=2$；(2) 定点为 $(0,1)$。设 $P(x_0,y_0)$，切线 $y=\dfrac{x_0}2x-y_0$ 与 $y=-1$ 交于 $Q$，可验证 $(0,1)$ 满足以 $PQ$ 为直径的圆的方程
  \item[\textbf{22.}] (1) $f(x)=x\sin x-\dfrac32$。由 $f'(x)=a(\sin x+x\cos x)$ 知 $f$ 在 $\left[0,\dfrac\pi2\right]$ 上单调，最大值只能在端点取得，比较得 $a=1$；(2) $2$ 个。$f$ 在 $\left(0,\dfrac\pi2\right)$ 上单调递增且 $f(0^+)<0<f\left(\dfrac\pi2\right)$，在 $\left(\dfrac\pi2,\pi\right)$ 上先增后减且 $f\left(\dfrac\pi2\right)>0>f(\pi^-)$，各恰有一个零点
\end{itemize}
